\section{Pilotphase der Studie}
Bevor die Studierenden mit der Studie konfrontiert werden sollten, wurden die Aufgaben und Abläufe in der Pilotphase erprobt. Die Pilotphase wurde mit insgesamt 3 Personen druchgeführt. Die Auswahl erfolgte aufgrund ihrer Eigenschaften, welche Überschneidungen im Bereich Alter, Studienrichtung oder Vorwissen im Vergleich zur Studiengruppe aufweisen. Im Folgenden sind die Eckdaten der Personen aufgelistet, die die Pilotphase durchlaufen sind:

Person A: 30 Jahre, männlich, Bachelorstudent im Bereich Cloud und IT. Programmierkenntnisse: Gerning. Erfahrung mit Python: Gering

Person B: 28 Jahre, männlich, Bachelorstudent für Lehramt Informatik. Programmierkenntnisse: Mittel. Erfahrung mit Python: Gering

Person C: 26 Jahre, weiblich, Lehrabschluss IT-Technik, Programmierkenntnisse: Gering. Erfahrung mit Python: Gering

Zur Einarbeitung der Verbesserungen wurde ein inkrementelle bzw. agiler Ansatz gewählt. So wurden die drei Personen auch nicht an einem Tag zur Pilotstudie eingeladen. Die Termine wurden mit den Personen nacheinander arrangiert, damit zwischen den Durchführungen genug Zeit vorhanden war, um die Studie anzupassen. Um diese inkrementellen Anpassungen soll es im nun folgenden Kapitel gehen. Es soll auf die Teilaspekte Aufgabenstellung, Entwicklungsumgebung, Quellcode und Bewertungsraster eingeangen weden.

\subsection{Durchführung und Erkenntnisse der Pilotstudie}
Die Termine mit den Personen in der Pilotphase wurden online über Discord abgehalten. Die Versuchspersonen wurden von Anfang an gebeten den Bildschirm freizugeben, damit auch die Schritte vor der Bearbeitung der Aufgaben nachvollzogen weden können.

\subsubsection{Durchführung mit Person A}
Vor dem Termin mit Person A entstand eine rudimentäre Checkliste, die den allgemeinen Ablauf der Studie festhaltet. Die Aufgabenstellung bestand derzeit aus vier Aufgaben mit dem Aufgabenprofil (Aufgabennummer - Klassifikation nach UMC-Framework - Schwierigkeitsgrad): Aufgabe 1 - use - einfach, Aufgabe 2 - modify - einfach, Aufgabe 3 - modify - einfach/mittel, Aufgabe 4 - create/modify - schwer. Person A hat die Aufgaben lokal am persönlichen Windows-11-PC in der Entwicklungsumgebung VS Code bearbeitet. Python war bereits installiert. 

Ein zentrales Thema, das während der Durchführung mit Person A auffiel, war die Entwicklungsumgebung. In der Hauptstudie galt das \textit{Bring Your Own Device} - Prinzip und somit musste sichergestellt werden, dass alle TeilnehmerInnen dieselbe Ausgangslage haben. Zwei Lösungsmöglichkeiten wurden im Anschluss diskutiert. Einerseits die Bereitstellung eines vorkonfigurierten Abbilds einer virtuellen Maschine, das mit allen zur Bearbeitung notwendigen Werkzeugen an die Studierenden verteilt werden kann. Andererseits wurde die Verwendung von GitHub Codespaces in Betracht gezogen, das ebenfalls zentral konfiguriert werden kann, jedoch eine cloudbasierte Lösung darstellt. Ein Codespace wird über den Browser aus einem GitHub-Repository heraus gestartet und man erhält in einem neuen Tab eine VS-Code-Webclient-Oberfläche. Die Entwicklungsumgebung läuft in einem Entwicklungscontainer, genannt DevContainer. Dieser wird über die Datei \textit{devcontainer.json}, die sich im Ordner \textit{.devcontainer} im Repository befindet, konfiguriert. Als Konfigurationsmöglickeiten können unter anderem die benötigten Werkzeuge, Laufzeitumgebungen sowie VS-Code-Erweiterungen und Einstellungen festgelegt werden.

ZIT: https://docs.github.com/en/codespaces/setting-up-your-project-for-codespaces/adding-a-dev-container-configuration/introduction-to-dev-containers




Zusammengefasst fanden in dieser anfänglichen Phase des Pilots die Erstellung der Checkliste sowie die erstmalige Evaluierung des Bewertungsrasters statt. Zudem wurde die zentrale Frage nach der Programmierumgebung aufgeworfen.

\subsubsection{Durchführung mit Person B}
Einen Tag nach Person A wurde mit Person B der Pilot durchgeführt. Die größte Änderung im Vergleich zu



\subsubsection{Durchführung mit Person C}





\subsection{Adaptionen für die Hauptstudie}

